\section{A proved SOR extension to every fixed bipartite face}
\label{sec:fixed-face-sor}

The exact star proof uses a two-cell bipartition.  On a general bipartite
face, the same algebra survives mode by mode after a singular-value
decomposition.  This yields a genuine extension of the accelerated spectral
mechanism, with an important qualifier: the face is already exposed and is
held fixed.

Let $U=R\mathbin{\dot\cup}B$ induce a bipartite subgraph.  The shared matrix
restricted to $U$ is
\begin{equation}
 Q_{UU}=\frac{1+\alpha}{2}
 \left(I-c_\alpha S_{UU}\right),
 \qquad
 c_\alpha:=\frac{1-\alpha}{1+\alpha},
 \qquad
 S:=D^{-1/2}AD^{-1/2}.
 \label{eq:scaled-fixed-face}
\end{equation}
The degrees in $D$ remain ambient graph degrees.  In the red--blue order,
write
\begin{equation}
 S_{UU}=\begin{bmatrix}0&C\\C^T&0\end{bmatrix},
 \qquad
 \sigma_U:=\|C\|_2,
 \qquad
 \rho_U:=c_\alpha\sigma_U.
 \label{eq:bipartite-cross-block}
\end{equation}
Since $S$ has spectrum in $[-1,1]$, $0\leq\rho_U<1$.

One red--black SOR sweep first relaxes every coordinate in $R$ and then every
coordinate in $B$, using a common $\omega\in(0,2)$.  Coordinates inside one
color have no edges, so their updates commute and constitute one exact block
relaxation.

\begin{theorem}[Exact optimal red--black SOR factor]
\label{thm:fixed-bipartite-sor}
Suppose $\rho_U>0$ and put
\begin{equation}
 t_U:=\sqrt{1-\rho_U^2},
 \qquad
 \omega_U^\star:=\frac{2}{1+t_U},
 \qquad
 \zeta_U:=\omega_U^\star-1
 =\frac{1-t_U}{1+t_U}.
 \label{eq:fixed-face-optimal-sor}
\end{equation}
For the scaled system $(I-c_\alpha S_{UU})x=h$, the red--black SOR error
matrix $T_U$ at $\omega_U^\star$ satisfies
\begin{equation}
 \rho(T_U)=\zeta_U.
 \label{eq:fixed-face-spectral-radius}
\end{equation}
More precisely, for every integer $k\geq1$,
\begin{equation}
 \|T_U^k\|_2\leq6k\zeta_U^{k-1}.
 \label{eq:fixed-face-power-bound}
\end{equation}
When $\rho_U=0$, one sweep with $\omega=1$ solves the decoupled system.
\end{theorem}

\begin{proof}
Take a singular pair of $C$ with singular value $\sigma$.  In its two scalar
amplitudes $(a,b)$, the homogeneous error update is
\begin{align}
 a^+&=(1-\omega)a+\omega c_\alpha\sigma b,
 \notag\\
 b^+&=(1-\omega)b+\omega c_\alpha\sigma a^+.
 \label{eq:sor-singular-mode-update}
\end{align}
Thus the mode matrix is
\begin{equation}
 T_\sigma(\omega)=
 \begin{bmatrix}
 1-\omega&\omega c_\alpha\sigma\\
 \omega c_\alpha\sigma(1-\omega)&
 (1-\omega)+\omega^2c_\alpha^2\sigma^2
 \end{bmatrix}.
 \label{eq:sor-mode-matrix}
\end{equation}
It has determinant $(1-\omega)^2$ and trace
\begin{equation}
 2(1-\omega)+\omega^2c_\alpha^2\sigma^2.
 \label{eq:sor-mode-trace}
\end{equation}

At $\omega=\omega_U^\star$, one has
$1-\omega=-\zeta_U$ and
\begin{equation}
 (\omega_U^\star)^2\rho_U^2=4\zeta_U.
 \label{eq:sor-optimal-identity}
\end{equation}
As $0\leq c_\alpha\sigma\leq\rho_U$, the trace in
\eqref{eq:sor-mode-trace} ranges from $-2\zeta_U$ to $2\zeta_U$, while the
determinant is $\zeta_U^2$.  The two eigenvalues therefore either form a
complex-conjugate pair of modulus $\zeta_U$ or meet at
$\pm\zeta_U$.  The largest singular mode meets at $+\zeta_U$.
Null singular directions have the scalar factor $-\zeta_U$.  Orthogonal
singular-vector changes of basis block-diagonalize the full error map, proving
\eqref{eq:fixed-face-spectral-radius}.

It remains to justify optimality.  On the largest singular mode the
characteristic discriminant factors as
\begin{equation*}
 \omega^2\rho_U^2
 \left(4(1-\omega)+\omega^2\rho_U^2\right).
\end{equation*}
The nonzero factor first vanishes in $(1,2)$ exactly at
$\omega=2/(1+t_U)$.  Before this double-root point the larger real root
decreases; after it, the determinant lower bound
$\rho(T_\sigma)\geq|1-\omega|$ increases.  Values $\omega\leq1$ have a
larger top-mode root.  Hence no common relaxation parameter has a smaller
worst-modal spectral radius.

For the finite bound, the entries of every $T_\sigma(\omega_U^\star)$ are
bounded in magnitude by constants and
$\|T_\sigma\|_2\leq5$.  If its eigenvalues are $\lambda_1,\lambda_2$, the
Cayley--Hamilton identity gives
\begin{equation*}
 T_\sigma^k=s_kT_\sigma-\zeta_U^2s_{k-1}I,
 \qquad
 s_k:=\sum_{j=0}^{k-1}\lambda_1^{k-1-j}\lambda_2^j.
\end{equation*}
Here $s_0:=0$.  The same formula holds by continuity at a repeated root, and
$|s_k|\leq k\zeta_U^{k-1}$.  Hence
$\|T_\sigma^k\|_2\leq6k\zeta_U^{k-1}$.  Taking the maximum over the
orthogonal mode blocks proves \eqref{eq:fixed-face-power-bound}.
\end{proof}

The factor is accelerated because
\begin{equation}
 \zeta_U=\frac{1-t_U}{1+t_U}\leq e^{-2t_U},
 \qquad
 t_U\geq\sqrt{1-c_\alpha^2}
 =\frac{2\sqrt\alpha}{1+\alpha}.
 \label{eq:fixed-face-universal-t}
\end{equation}
The first inequality is elementary, and the second uses $\sigma_U\leq1$.

\begin{corollary}[Oracle-tuned finite fixed-face semantic solve]
\label{cor:fixed-face-semantic-sor}
Let $x^\dagger\geq0$ solve either the unregularized problem on $U$ or the
correct positive RPPR face system \eqref{eq:rppr-active-face-system}, and
suppose $\|x^\dagger\|_2\leq1$ and $0<\tau\leq1$.  If
$\sigma_U$ (equivalently $\omega_U^\star$) is supplied, starting from zero
with that parameter produces $x^{(k)}$ with
\begin{equation}
 \|x^{(k)}-x^\dagger\|_{D,\infty}\leq\tau
 \label{eq:fixed-face-semantic-target}
\end{equation}
after
\begin{equation}
 k=O\!\left(
  \frac1{t_U}\log\frac{2}{t_U\tau}
 \right)
 \label{eq:fixed-face-sweep-count}
\end{equation}
full red--black sweeps.  Including repeated row scans and output, the charged
work is
\begin{equation}
 W_U=O\!\left(
  \frac{\operatorname{vol}(U)}{t_U}
  \log\frac{2}{t_U\tau}
 \right)
 =O\!\left(
  \frac{\operatorname{vol}(U)}{\sqrt\alpha}
  \log\frac{2}{\tau\sqrt\alpha}
 \right).
 \label{eq:fixed-face-work}
\end{equation}
\end{corollary}

\begin{proof}
All degrees are at least one, so
$\|z\|_{D,\infty}\leq\|z\|_2$.  Apply
\eqref{eq:fixed-face-power-bound} to the zero-start error and use
$\zeta_U\leq e^{-2t_U}$.  If $t_U\leq1/2$, the polynomial factor $k$ is
absorbed by choosing a sufficiently large absolute constant in
\eqref{eq:fixed-face-sweep-count}; if $t_U>1/2$, geometric convergence needs
only $O(\log(2/\tau))$ sweeps.  Each sweep scans every row in $U$ once and
costs $\operatorname{vol}(U)$.  Since $|U|\leq\operatorname{vol}(U)$, the
explicit output is absorbed.  Finally use \eqref{eq:fixed-face-universal-t}.
\end{proof}

\begin{proposition}[Certified face tuning with spectral cost]
\label{prop:certified-face-spectral-cost}
Suppose a charged spectral routine returns a certified upper bound
\begin{equation}
 \rho_U\le \bar\rho_U<1,
 \qquad
 \bar t_U:=\sqrt{1-\bar\rho_U^2},
 \qquad
 \bar\omega_U:=\frac2{1+\bar t_U},
 \qquad
 \bar\zeta_U:=\bar\omega_U-1.
 \label{eq:certified-face-advice}
\end{equation}
Then red--black SOR with the scalar $\bar\omega_U$ has
\begin{equation}
 \rho(T_U)\le\bar\zeta_U,
 \qquad
 \|T_U^k\|_2\le6k\bar\zeta_U^{k-1}.
 \label{eq:certified-face-power}
\end{equation}
If producing the certificate costs $\mathsf C_{\rm spec}(U)$ and, for some
$0<\gamma\le1$, it also certifies $\bar t_U\ge\gamma t_U$, the fully charged
semantic solve costs
\begin{equation}
 \boxed{
 \mathsf C_{\rm spec}(U)
 +O\!\left(
  \frac{\operatorname{vol}(U)}{\gamma t_U}
  \log\frac{2}{\gamma t_U\tau}
 \right).}
 \label{eq:certified-face-total-work}
\end{equation}
The relative-gap guarantee follows, in particular, from the squared-radius
certificate
\begin{equation}
 0\le\bar\rho_U^2-\rho_U^2
 \le(1-\gamma^2)t_U^2.
 \label{eq:certified-face-required-accuracy}
\end{equation}
Thus, within this certified-upper tuning rule, a constant-factor realization
of the face-optimal sweep count requires additive accuracy on the scale
$t_U^2$ in the squared Jacobi radius.  This scale can be as small as
$\Theta(\alpha)$.
\end{proposition}

\begin{proof}
At $\bar\omega_U$, one has
$1-\bar\omega_U=-\bar\zeta_U$ and
$\bar\omega_U^2\bar\rho_U^2=4\bar\zeta_U$.  Every true modal squared coupling
is at most $\rho_U^2\le\bar\rho_U^2$, so its trace lies in
$[-2\bar\zeta_U,2\bar\zeta_U]$ and its determinant is
$\bar\zeta_U^2$.  The singular-mode and Cayley--Hamilton proof of
\Cref{thm:fixed-bipartite-sor} gives
\eqref{eq:certified-face-power}.  Applying that bound as in
\Cref{cor:fixed-face-semantic-sor} and adding the spectral routine's work
proves \eqref{eq:certified-face-total-work}.  Finally,
\eqref{eq:certified-face-required-accuracy} gives
\[
 \bar t_U^2=1-\bar\rho_U^2
 \ge 1-\rho_U^2-(1-\gamma^2)t_U^2
 =\gamma^2t_U^2.
\]
The universal lower bound $t_U^2\ge1-c_\alpha^2=\Theta(\alpha)$ gives the
last statement.

The upper-certificate direction is essential.  If a lower estimate
$\widetilde\rho_U<\rho_U$ is inserted into the same formula, its tuned top-mode
trace is
\[
 2\widetilde\zeta_U
 +\widetilde\omega_U^2
   (\rho_U^2-\widetilde\rho_U^2)
 >2\widetilde\zeta_U,
\]
while the determinant remains $\widetilde\zeta_U^2$.  The larger real root is
then strictly greater than the advertised factor.  A raw Rayleigh quotient
cannot silently serve as safe face advice.
\end{proof}

This proposition makes the formerly hidden oracle cost explicit.  On one
supplied face, materializing the cross block costs
$O(\operatorname{vol}(U))$.  A dense SVD to the required certified precision
uses $O(|U|^3)$ real-arithmetic work and $O(|U|^2)$ workspace, so it gives the
valid but generally nonlocal bound
\begin{equation}
 O\!\left(
  \operatorname{vol}(U)+|U|^3
  +\frac{\operatorname{vol}(U)}{t_U}
   \log\frac{2}{t_U\tau}
 \right).
 \label{eq:dense-face-spectral-work}
\end{equation}
The bit cost needed to certify a requested eigenvalue interval is additional
and depends on the input encoding.  Alternatively, an iterative estimator
using $s$ sparse products costs $O(s\operatorname{vol}(U))$ plus vector
operations.  The value of $s$ is not graph-uniform without an estimator
theorem: ordinary power/Rayleigh estimates approach $\rho_U$ from below and
are not the upper certificate required in
\eqref{eq:certified-face-advice}; their rate also depends on the top singular
gap.  On changing faces, every recomputation is charged, so the spectral term
is $\sum_j\mathsf C_{\rm spec}(U_j)$ rather than one free scalar.

The qualifier ``supplied'' therefore matters: the sharp $t_U$ bound does not
include a local computation of $\sigma_U$.  Without spectral advice,
\Cref{lem:one-omega-all-bipartite-faces} uses the graph-global parameter
$\omega_\alpha$ and gives the fully charged generic
$O(\operatorname{vol}(U)\alpha^{-1/2}
\log(2/(\tau\sqrt\alpha)))$ bound without an eigenvalue computation.

The norm assumption is automatic for the shared single-seed or probability-
seed PPR problem.  Indeed,
\begin{equation}
 \|b\|_2=\alpha\|D^{-1/2}s\|_2\leq\alpha,
 \qquad
 \|x^0\|_2\leq\|Q^{-1}\|_2\|b\|_2\leq1.
 \label{eq:ppr-unit-two-norm}
\end{equation}
Moreover, \Cref{thm:rppr-ppr-bias} gives
$0\leq x^\rho\leq x^0$, so $\|x^\rho\|_2\leq1$ as well.

\subsection{What has and has not generalized}

The theorem generalizes three star ingredients:

\begin{itemize}[leftmargin=*]
\item the relaxation parameter is determined by an exact Jacobi spectral
      radius;
\item every seed-visible mode is accelerated at the factor $\zeta_U$;
\item actual degree work, rather than only iteration count, is charged for a
      fixed face.
\end{itemize}

It does not generalize the star's log-free semantic cancellation.  The safe
finite bound carries a logarithm and an absolute norm conversion.  It also
does not find $U$, verify exterior KKT signs, transport momentum to a larger
face, or compute $\sigma_U$ locally.  These are the next layers of the graph
problem, not details hidden inside the spectral theorem.

For a known RPPR support $A$, the source volume bound
$\operatorname{vol}(A)\leq1/\rho$ turns
\eqref{eq:fixed-face-work} into
\begin{equation}
 W_A=O\!\left(
  \frac1{\rho\sqrt\alpha}
  \log\frac{2}{\tau\sqrt\alpha}
 \right).
 \label{eq:oracle-face-rppr-work}
\end{equation}
This is a useful oracle benchmark.  Calling it an end-to-end local algorithm
would be incorrect because the support and its boundary certificate were
supplied for free.
